\documentclass[11pt]{article}
\usepackage[a4paper, top=18mm, bottom=18mm, left=20mm, right=20mm]{geometry}
\usepackage[T2A]{fontenc}
\usepackage[utf8]{inputenc}
\usepackage[ukrainian]{babel}
\usepackage{graphicx}
\usepackage[export]{adjustbox}
\usepackage{amsmath}
\usepackage{amssymb}
\graphicspath{{/app/Quiz/Scripts/2023/ProgAll/15BinTree/}}
\setlength{\parindent}{0pt}
\setlength{\parskip}{6pt}
\emergencystretch=3em
\pagestyle{empty}
\begin{document}
%begin
{\centering\Large\bfseries Контрольна робота\par}
\medskip
\noindent\rule{\linewidth}{0.5pt}
\smallskip
\noindent\textbf{Варіант \No\,4}\hfill \textit{ПІБ:}\,\underline{\hspace{60mm}}
\vspace{4mm}

%beginitem
1. \textbf{Що буде виведено за виконання фрагменту коду?}

print('R', end=' ')
for e in range(10,10,-3):
    if e>=7:
        break
        print(e, end=' ')
    if e>=7:
        break
    else:
        print(e, end=' ')
    print(e)
\vspace{5mm}
%enditem

%beginitem
2. \textbf{Що буде виведено за виконання фрагменту коду?}
a = 4
b = 2
c = 1

def h(a):
global c    a = 3
    b *= 4
    c = 5
    return a + b + c

a = 5
b = 6
c = 0

try:
    print(h(b), a, b, c, sep=';')
except:
    print('Z', a, b, c, sep=';')
\vspace{5mm}
%enditem

%beginitem
3. 

\begin{center}
	\adjustimage{max size={0.8\linewidth}{0.8\paperheight}}
	{../15BinTree/trees_exp/133.png}
\end{center}
Указати послідовність міток вершин дерева, зображеного на рис., що утвориться в результаті обходу дерева в прямому порядку. Мітки записувати через пробіл.\par Алгоритм починає роботу з кореня (на рис. це верхня вершина).
%answer:133 прямому
\vspace{5mm}
%enditem

%beginitem
4. 
Записати ВИРАЗ, значенням якого є список, що складається з кубів цілих чисел від 15 до 2005 (включно), що не діляться на жодне з чисел 2, 3, записаних в оберненому порядку.\vspace{5mm}
%enditem

%beginitem
5. %goodpath:Scripts/2019/Mod2019_1/Lex/01-good.txt
\textbf{Відмітити все, що є однією правильною лексемою:}
( 1 ) -3J

( 2 ) "a"a"

( 3 ) 0.5

( 4 ) 0b001
\vspace{5mm}
%enditem

%beginitem
6. 
Написати програму, що запитує у користувача значення дійсних параметрів $x$ та $y$, обчислює для них значення виразу 
$$\frac{\log_{2} x + 76}{(\arccos x + 0.6) (y + 87)}$$ 
та повiдомляє введенi данi та результат обчислення.
 Оцінюється правильність та якість коду.

\vspace{5mm}
%enditem

%beginitem
7. \textbf{Що буде виведено за виконання фрагменту коду?}
e=[10,11,12,13,14]; e.extend([1,2]); print(e)\vspace{5mm}
%enditem

%beginitem
8. 
Записати ВИРАЗ, значенням якого є множина, що складається з цілих чисел від 15 до 2005 (включно), що не діляться на 7.\vspace{5mm}
%enditem

%beginitem
9. %goodpath:Scripts/2018/Mod2018_1/01-good.txt
\textbf{Відмітити все, що є однією правильною лексемою:}
( 1 ) x*y

( 2 ) (1)

( 3 ) False

( 4 ) For
\vspace{5mm}
%enditem

%beginitem
10. %goodpath:Scripts/2019/Mod2019_1/Lex/03-good.txt
\textbf{Відмітити все, що є коректним оператором С++:}
( 1 ) (2+2j)//2

( 2 ) if a<b: a=0 else: b=0

( 3 ) x=y=z

( 4 ) x==y==z
\vspace{5mm}
%enditem

%end
\newpage
\end{document}

